\usepackage{fontspec}
\setmainfont{texgyretermes-regular.otf}[BoldFont=texgyretermes-bold.otf,ItalicFont=texgyretermes-italic.otf]
\usepackage[polish]{babel}
\usepackage{amsmath,amssymb,geometry,array,tabularx,graphicx,caption,fancyhdr}
\usepackage[table]{xcolor}
\geometry{left=18mm,right=18mm,top=19mm,bottom=19mm,headheight=14pt,headsep=6mm}
\setlength{\parindent}{0pt}\setlength{\parskip}{5pt}
\setlength{\emergencystretch}{2em}
\renewcommand{\arraystretch}{1.3}\setlength{\tabcolsep}{5pt}
\setlength{\fboxsep}{6pt}
\setlength{\abovedisplayskip}{8pt}\setlength{\belowdisplayskip}{8pt}
\setlength{\abovedisplayshortskip}{5pt}\setlength{\belowdisplayshortskip}{5pt}
\widowpenalty=10000\clubpenalty=10000
\raggedbottom
\newcolumntype{Y}{>{\raggedright\arraybackslash}X}
\newcolumntype{C}{>{\centering\arraybackslash}X}
\newcommand{\naglowek}{\rowcolor{gray!12}}
\newcommand{\odstep}{\par\smallskip}
\captionsetup{font=small,labelfont=bf,justification=justified,singlelinecheck=false}
\pagestyle{fancy}\fancyhf{}\fancyhead[L]{\small Ćwiczenie 31 — modelowanie pola}
\fancyhead[R]{\small\Rodzaj}\fancyfoot[C]{\small\thepage}
\renewcommand{\headrulewidth}{0.3pt}
\makeatletter
\renewcommand\section{\@startsection{section}{1}{0pt}{10pt}{4pt}{\large\bfseries}}
\renewcommand\subsection{\@startsection{subsection}{2}{0pt}{8pt}{4pt}{\normalsize\bfseries}}
\makeatother
\newcommand{\pole}{\rule{23mm}{0.3pt}}
\newcommand{\miejsce}[2]{\par\medskip\noindent\fbox{\begin{minipage}[t][#1][t]{\dimexpr\linewidth-2\fboxsep-2\fboxrule\relax}\small #2\end{minipage}}\par\medskip}
\newcommand{\metryka}{%
\begingroup\small\setlength{\parskip}{0pt}\renewcommand{\arraystretch}{1.1}
\begin{tabularx}{\linewidth}{|p{29mm}|Y|p{18mm}|p{18mm}|p{19mm}|}\hline
\textbf{Wydział}\newline Informatyki & \textbf{Imię i nazwisko}\newline 1. Marcin Podyma\newline 2. Paweł Mierzwa & \textbf{Rok}\newline 2026 & \textbf{Grupa}\newline 3 & \textbf{Zespół}\newline 1\\\hline
\end{tabularx}\par\nointerlineskip
\begin{tabularx}{\linewidth}{|p{29mm}|Y|p{19mm}|}
\parbox[t]{29mm}{\centering\textbf{PRACOWNIA\\FIZYCZNA\\WFiIS AGH}} & \textbf{Temat:}\newline Modelowanie pola elektrycznego & \textbf{Nr}\newline\textbf{ćwiczenia}\newline 31\\\hline
\end{tabularx}\par\nointerlineskip
\begin{tabularx}{\linewidth}{|p{29mm}|Y|Y|Y|Y|p{19mm}|}
\textbf{Data wykonania}\newline\strut & \textbf{Data oddania}\newline\strut & \textbf{Zwrot do popr.} & \textbf{Data oddania} & \textbf{Data zaliczenia} & \textbf{OCENA}\\\hline
\end{tabularx}\endgroup\medskip}


